\documentstyle[%


righttag%


,11pt%


,amssymb%


,russian%               remove in Eng.


,izvru%                 Set izven% in Eng.


% ,izvru-ks             for brief communications


]{amsart}





% 





\newcommand{\la}{\langle}


\newcommand{\ra}{\rangle}


\newcommand{\tri}{|@!@!@!@!@!|@!@!@!@!@!|}


%% \newcommand{\Cl}{\operatorname{Cl}}  


\newcommand{\Cl}{\cal C l}


\newcommand{\im}{\operatorname{im}}


\newcommand{\coim}{\operatorname{coim}}


\newcommand{\dom}{\operatorname{dom}}


\newcommand{\const}{\operatorname{const}}


\newcommand{\tts}[1]{}


\numberwithin{equation}{section}





\theoremstyle{plain}


\theorembodyfont{\it}


\newtheorem{thms}{\hskip\parindent Теорема}[section]


\newtheorem{lems}{\hskip\parindent Лемма}[section]





\theoremstyle{definition}


\theorembodyfont{\rm}


\newtheorem{defns}{\hskip\parindent Определение}[section]


\newtheorem{notes}{\hskip\parindent Замечание}[section]


\newtheorem{cors}{\hskip\parindent Следствие}[section]


% \renewcommand{\thecornonum}{}











%\hoffset=-15mm%                  For Eng.


\setcounter{page}{54}


\god{2005}


\nomer{10~(521)} %                        Remove in Eng.


%\nomer{Vol.\,49, No.\,10}%               Set in Eng.


%\str{\pageref{begin}--\pageref{end}}%  Set in Eng.


\udk{517.956}





\author{\it Г.А.\,Свиридюк, И.К.\,Тринеева}


% NB:   ^^ remove in Eng.


\title{Сборка Уитни в фазовом


пространстве уравнения Хоффа}


\thanks{Работа выполнена при частичной финансовой поддержке


Министерства образования и науки России, грант \No\,РО~02-1.1-82.}





\date{}


%\pagestyle{myheadings}%         Set in Eng.


\pagestyle{plain} %              Remove in Eng.


\begin{document}


%\thispagestyle{plain}%          Set in Eng.


\maketitle


\label{begin}





\setcounter{section}{0}





Пусть $\Omega\subset{\Bbb R}^n$ --- ограниченная область с границей


$\partial \Omega$ класса $C^{\infty}$. В цилиндре


$\Omega\times\Bbb R$ уравнение Хоффа


\begin{equation}


(\lambda+\Delta)u_t=\alpha


u+\beta u^{3}


\end{equation}


при $n=1$ моделирует динамику


выпучивания двутавровой балки [1]. Здесь параметр $\lambda\in\Bbb


R_{+}$ соответствует внешней нагрузке, а параметры


$\alpha,\beta\in\Bbb R$ характеризуют свойства балки; искомая


функция $u=u(x,t)$ определяет отклонение балки от вертикали.





Задача Коши--Дирихле


\begin{equation}


\begin{split}


u(x,0)&=u_{0}(x),\quad \ \


x\in\Omega;\\


u(x,t)&=0,\qquad (x,t)\in\partial\Omega\times\Bbb R,


\end{split}


\end{equation}


для уравнения (0.1) изучалась в


[2] методами теории ветвления. В [3] начато изучение фазового


пространства уравнения (0.1) при $1\leq n\leq 4$. Там показано,


что в фазовом пространстве уравнения Хоффа лежат те точки $u\in


L_{4}(\Omega)$, для которых


\begin{equation}


\int_{\Omega}(\alpha + 3\beta u^{2})u


\varphi_{l} dx=0, \quad\lambda=-\lambda_{l},


\end{equation}


причем определитель


\begin{equation}


\bigg|\int_{\Omega}(\alpha + 3 \beta


u^{2})\varphi_{k} \varphi_{l} dx\bigg|\neq 0,


\quad \lambda=-\lambda_{l},\ \  \lambda=-\lambda_{k}.


\end{equation}


Здесь $\{\varphi_{m}\}$ --- ортонормированное (в смысле


$L_{2}(\Omega)$) семейство собственных функций однородной задачи


Дирихле для оператора Лапласа $\Delta$ в области $\Omega$,


соответствующих собственным значениям $\{\lambda_{m}\}$, которые


занумерованы по невозрастанию с учетом кратности. (Если


$-\lambda\notin\{\lambda_{m}\}$, то фазовое пространство


совпадает с пространством $L_{4}(\Omega)$, а условие (0.4)


исчезает.)





Другими словами, в [3] показано, что фазовое пространство (0.3)


уравнения Хоффа локально (т.\,е. в некоторых окрестностях тех своих


точек, для которых выполняется (0.4)) является банаховым


$C^{\infty}$-многообразием, моделируемым некоторым


дополнительным к ядру $\ker(\lambda + \Delta)$


подпространством. В [4] показано, что при $\alpha\beta>0$


фазовым пространством (0.3) служит простое банахово


$C^{\infty}$-многообразие (т.\,е. любой его атлас эквивалентен


атласу, содержащему единственную карту), причем в любой его точке


выполняется (0.4).





Исследуем фазовое пространство  уравнения Хоффа в


случае $\alpha\beta<0$ и при нарушении условия (0.4). Чтобы


привлечь идеи и методы [5], ограничимся случаем $\dim


\ker(\lambda + \Delta)=1$, который заведомо имеет место при $n=1$.


Основной результат статьи --- характеристика особенностей фазового


пространства (0.3) посредством относительно присоединенных


векторов.





Первый параграф данной работы носит вспомогательный характер, он содержит


результаты из [6], представленные в удобном для дальнейшего виде.


Во втором, опираясь на результаты [3], [7],


рассматриваем задачу (0.1), (0.2) в абстрактном виде. В


третьем параграфе содержится основной результат статьи. Список


литературы не претендует на полноту.





Исследования проводятся в


вещественных банаховых пространствах, но в 


``спектральных'' вопросах вводится их естественная комплексификация.


Все контуры ориентированы движением против часовой стрелки и


ограничивают области, лежащие слева при таком движении. Символами


$\Bbb I$ и $\Bbb O$ обозначены единичный и нулевой операторы,


области определения которых ясны из текста.








\section{Относительно $\sigma$-ограниченные операторы\protect \\


и относительно присоединенные векторы}





Пусть $\frak U$ и $\frak F$ --- банаховы пространства, оператор


$L\in \Cl(\frak U;\frak F)$ (т.\,е. линеен, замкнут и плотно


определен), а оператор $M\in L(\frak U;\frak F)$ (т.\,е. линеен


и непрерывен). Введем $L$-резольвентное


множество $\rho^{L}(M)=\{\mu\in \Bbb C:(\mu L-M)^{-1}\in


L(\frak F;\frak U)\}$ и $L$-спектр $\sigma^{L}(M)=\Bbb


C\setminus\rho^{L}(M)$ оператора $M$. Нетрудно показать [6], что


$L$-резольвентное множество всегда открыто, поэтому $L$-спектр


всегда замкнут.








\begin{defns}


Оператор $M$ называется {\it спектрально


ограниченным относительно} оператора $L$


($(L,\sigma)$-{\it ограниченным}), если


$$


\exists a\in\Bbb R_{+}\ \; \forall\mu\in \Bbb C\ \; (|\mu|>a)\Rightarrow


(\mu\in\rho^{L}(M)).


$$


\end{defns}





Пусть оператор $M$ \ $(L,\sigma)$-ограничен. Выберем контур


$\Gamma=\{\mu\in \Bbb C:|\mu|=r>a\}$ и построим операторы


$$


P=\frac{1}{2\pi i}\int_{\Gamma}R^{L}_{\mu}(M)d\mu,\quad


Q=\frac{1}{2\pi i}\int_{\Gamma}L^{L}_{\mu}(M)d\mu.


$$


Здесь $R^{L}_{\mu}(M)=(\mu L-M)^{-1}L$ --- правая, а


$L^{L}_{\mu}(M)=L(\mu L-M)^{-1}$ --- левая $L$-резольвенты оператора


$M$. Нетрудно показать, что операторы $P\in L(\frak U)$ и


$Q\in L(\frak F)$ являются проекторами~[6].


Положим $\frak U^0=\ker P$, $\frak U^1=\im P$, $\frak F^0=\ker Q$,


$\frak F^1=\im Q$, а через $L_{k}$ ($M_{k}$) обозначим сужение


оператора $L(M)$ на $\dom L\cap\frak U^{k}$ ($\frak


U^{k}$), $k=0,1$.








\begin{thms}


Пусть оператор $M$ \ $(L,\sigma)$-ограничен. Тогда


\begin{itemize}


\item[(i)] операторы $L_{k}\in\Cl(\frak U^{k};\frak F^{k}),


M_{k}\in L(\frak U^{k};\frak F^{k})$, $k=0,1;$


\item[(ii)] существуют операторы $L^{-1}_{1}\in L(\frak F^{1}; \frak


U^{1})$, $M^{-1}_{0}\in L(\frak F^{0};\frak U^{0})$.


\end{itemize}


\end{thms}





Линеал $\frak U^{0}_{0}=\dom L\cap\frak U^{0}$,


снабженный ``нормой графика''


$\tri\cdot\tri=\|\cdot\|_{\frak U}+\|L\cdot\|_{\frak F}$,


является банаховым пространством


$\frak U^{0}_{0}=(\dom L\cap\frak


U^{0},\tri\cdot\tri)$, плотно и непрерывно вложенным в пространство


$\frak U^{0}$. Предположим, что в условиях теоремы 1.1 оператор


\begin{equation}


H=M^{-1}_{0}L_{0}\in L(\frak U^{0}_{0}).


\end{equation}





Очевидно, оператор $S=L^{-1}_{1}M_{1}\in L(\frak U^{1})$.








\begin{cors}


Пусть оператор $M$ \ $(L,\sigma)$-ограничен


и выполнено условие (1.1). Тогда


$$


(\mu L-M)^{-1}=-\sum_{q=0}^{\infty}\mu^{q}H^{q}M^{-1}_{0}(\Bbb I -Q)+


\sum_{q=1}^{\infty}\mu^{q}S^{q-1}L^{-1}_{1}Q


$$


при всех $\mu\in\sigma^{L}(M)$.


\end{cors}








\begin{defns}


Для $L$-резольвенты $(\mu L-M)^{-1}$


оператора $M$ точка $\infty$ называется


\begin{itemize}


\item[(i)] {\it устранимой особой точкой}, если $H\equiv\Bbb O$;


\item[(ii)] {\it полюсом порядка} $p\in\Bbb N$, если $H^{p}\neq\Bbb O$,


а


$H^{p+1}\equiv\Bbb O$;


\item[(iii)] {\it существенно особой точкой}, если $H^{q}\neq\Bbb O$ при


любом $q\in\Bbb N$.


\end{itemize}


\end{defns}





Пусть вектор $\varphi_{0}\in\ker L\setminus\{0\}$, упорядоченное


множество $\{\varphi_{1},\varphi_{2},\dotsc \}\subset\frak U$


называется {\it цепочкой $M$-присоединенных векторов} вектора


$\varphi_{0}$, если $L\varphi_{q+1}=M\varphi_{q}$,


$q\in\{0\}\cup\Bbb N$,


и


$\varphi_{q}\notin\ker


L\setminus\{0\}$, $q\in\Bbb N$. Цепочка может быть бесконечной, в


частности, она может быть заполнена нулями, если $\ker L\cap\ker


M\neq\{0\}$. Однако она обязательно конечна, если в ней существует


такой вектор $\varphi_{q}$, что $M\varphi_{k}\notin\im L$.


Мощность конечной цепочки называется  ее {\it длиной}. Порядковый


номер $M$-присоединенного вектора в цепочке (конечной или


бесконечной) называется его {\it высотой}. Линеал, натянутый на


все векторы ядра $\ker L$ и все $M$-присоединенные векторы,


называется {\it $M$-корневым линеалом} оператора $L$. Если 


$M$-корневой линеал замкнут, то он называется\linebreak $M$-{\it корневым


пространством} оператора $L$.








\begin{thms}


Пусть оператор $M$ \ $(L,\sigma)$-ограничен. Тогда 


\begin{itemize}


\item[(i)] $\frak U^{0}=\overline{\ker L}$ и любой вектор $\varphi\in\ker


L\setminus\{0\}$ не имеет $M$-присоединенных векторов,


если $\infty$ является устранимой особой точкой $L$-резольвенты


оператора $M;$


\item[(ii)] 


$\frak U^{0}$ является замыканием $M$-корневого


линеала оператора $L$ и длина любой цепочки\linebreak $M$-присоединенных


векторов ограничена числом $p$, если $\infty$ является


полюсом порядка $p\in\Bbb N$ \ $L$-резольвенты оператора $M;$


\item[(iii)] 


$\frak U^{0}$ содержит $M$-корневой линеал оператора $L$,


причем длины всех цепочек $M$-при\-со\-е\-ди\-нен\-ных векторов


неограничены, если $\infty$ является


существенно особой точкой\linebreak $L$-резольвенты  оператора $M$.


\end{itemize}


\end{thms}





В случае, когда оператор $L$ фредгольмов (т.\,е. $\mathrm{ind}\, L=0$),


существует некоторое обращение теоремы 1.2 [6]. Чтобы


сформулировать его, удобно считать устранимую особую точку


полюсом порядка нуль.








\begin{cors}


Пусть оператор $L$ фредгольмов. Тогда


следующие утверждения эквивалентны:


\begin{itemize}


\item[(i)] оператор $M$ \ $(L,\sigma)$-ограничен, причем $\infty$ --- полюс


порядка $p\in\{0\}\cup\Bbb N$ \ $L$-резольвенты оператора $M$;


\item[(ii)] длина любой цепочки $M$-присоединенных векторов оператора


$L$ ограничена числом\linebreak $p\in\{0\}\cup\Bbb N$.


\end{itemize}


\end{cors}








\section{Задача Коши для уравнения соболевского типа}





Пусть $\frak U$ и $\frak F$ --- банаховы пространства, операторы


$L\in \Cl(\frak U;\frak F)$, $M\in L(\frak U;\frak F)$, $N\in


C^{\infty}(\frak U;\frak F)$. Рассмотрим задачу Коши


\begin{equation}


u(0)=u_{0}


\end{equation}


для полулинейного уравнения соболевского типа


\begin{equation}


L\Dot{u}=Mu+N(u).


\end{equation}


Вектор-функцию $u\in


C^{\infty}((-T,T);\frak U)$ назовем {\it решением уравнения}


(2.2), если она удовлетворяет (2.2) при некотором 


$T\in\Bbb R_{+}$. Решение уравнения $u=u(t)$ называется


{\it решением задачи} (2.1), (2.2), если оно удовлетворяет (2.1).





Пусть оператор $M$ \ $(L,\sigma)$-ограничен, тогда в силу теоремы


1.1 уравнение (2.2) эквивалентно системе уравнений


\begin{gather}


L_{0}\Dot u{}^{0}=(\Bbb I-Q)(Mu+N(u)),\\


\Dot u{}^{1}=Su^{1}+F(u),


\end{gather}


где операторы


$S=L^{-1}_{1}M_{1}\in L(\frak U^{1})$, $F=L^{-1}_{1}QN\in


C^{\infty}(\frak U;\frak U^{1})$. 





Пусть выполнено условие


\begin{equation}


L_{0}\Dot u{}^{0}\equiv 0,


\end{equation}


тогда из (2.3) следует,


что любое решение уравнения (2.2) лежит (как траектория) в


множестве


$$


\frak H=\{u\in\frak U:(\Bbb I-Q)(Mu+N(u))=0\}.


$$


Такие решения названы {\it квазистационарными траекториями} [7].


Заметим, что условие (2.5) с необходимостью выполняется, если


нелинейный член $N\equiv\Bbb O$.





Пусть $u_{0}\in\frak H$. Немного отходя от стандарта [8],


некоторую окрестность $O_{u_{0}}\subset\frak H$ точки $u_{0}$


будем называть {\it картой на множестве} $\frak H$, если


существует такой $C^{\infty}$-диффеоморфизм\linebreak


$\delta:O^{1}_{u_{0}}\rightarrow O_{u_{0}}$ окрестности


$O^{1}_{u_{0}}\subset\frak U^{1}$ точки $u^{1}_{0}=Pu_{0}$,


что $\delta^{-1}=P$. Если такая карта существует, то множество


$\frak H$ будем называть {\it банаховым $C^{\infty}$-многообразием


в точке} $u_{0}$.








\begin{notes}


Обозначим через $N'_{u_{0}}$ производную


Фреше оператора $N$ в точке $u_{0}$, и пусть оператор


$M+N'_{u_{0}}:\frak U^{0}\rightarrow \frak F^{0}$ --- топлинейный


изоморфизм. Тогда множество $\frak H$ в точке $u_{0}$ является


банаховым $C^{\infty}$-многообразием. Действительно, в силу


теоремы о неявной функции существуют  окрестности


$O^{1}_{u_{0}}\subset\frak U^{1}$ и $O^{0}_{u_{0}}\subset\frak


U^{0}$ точек $u^{1}_{0}=Pu_{0}$ и $u^{0}_{0}=(\Bbb I-P)u_{0}$


соответственно и отображение $d\in


C^{\infty}(O^{1}_{u_{0}};O^{0}_{u_{0}})$ такие, что $\delta=\Bbb I


+d$ является требуемым $C^{\infty}$-диффеоморфизмом.


\end{notes}








\begin{thms}


Пусть оператор $M$ \ $(L,\sigma)$-ограничен


и выполнено условие $(2.5)$. Если в точке $u_{0}$ множество $\frak


H$ является $C^{\infty}$-многообразием, то для некоторого


$T\in\Bbb R_{+}$ существует единственное решение задачи $(2.1)$,


$(2.2)$.


\end{thms}








\begin{pf}


Пусть $\delta'_{u^{1}}:\frak


U^{1}\rightarrow T_{u}\frak H$ --- производная Фреше диффеоморфизма


$\delta$ в точке $u^{1}\in O^{1}_{u_{0}}$


($T_{u}\frak H$ ---


касательное пространство  в точке $u=\delta(u^{1})$).


Тогда из (2.4) вытекает


\begin{equation}


\Dot{u}=\delta'_{u^{1}}\Dot u{}^{1}=\delta'_{u^{1}}(SPu+F(u))=A(u),


\end{equation}


где $u\in O_{u_{0}}\subset\frak H$. В силу классической теоремы


Коши [8] для некоторого  $T\in\Bbb R_{+}$ существует


единственное решение $u\in C^{\infty}((-T,T);\frak H)$ задачи


(2.1), (2.6). Как нетрудно проверить, найденная вектор-функция


$u=u(t)$ является решением задачи (2.1), (2.2).


\end{pf}





Теорема 2.1 дает достаточные условия существования


квазистационарных траекторий. Условие (2.5) заведомо выполняется,


если $\ker L=\frak U^{0}$. Если вдобавок оператор $L$


фредгольмов, то условие $\ker L=\frak U^{0}$ эквивалентно


отсутствию $M$-присоединенных векторов (следствие 1.2). Если,


кроме того, отсутствуют $(M+N'_{u_{0}})$-присоединенные векторы


оператора $L$ в любой точке $u_{0}\in\frak H$, то это означает,


что оператор $(M+N'_{u_{0}}):\frak U_{0}\rightarrow\frak F^{0}$ ---


топлинейный изоморфизм, и в силу замечания 2.1 и теоремы 2.1 в


любой точке $u_{0}\in\frak H$ задача (2.1), (2.2) однозначно


разрешима. Именно эту ситуацию наблюдаем в задаче (0.1), (0.2)


в случае $\alpha\beta>0$ [4]. Рассмотрим теперь ситуацию,


когда фредгольмов оператор $L$ не имеет $M$-присоединенных


векторов, но имеет $(M+N'_{u_{0}})$-присоединенные. Для простоты


ограничимся случаем одномерного ядра $\ker L$.








\begin{thms}


Пусть оператор $L$ фредгольмов, $\dim\ker


L=1$. Пусть оператор $L$ не имеет  $M$-присоединенных векторов, но


имеет в некоторой 


точке $u_{0}\in\frak H$ \ $(M+N'_{u_{0}})$-присоединенные 


векторы высоты не большей $p\in\Bbb N$.


Тогда в некоторой окрестности $O_{u_{0}}\subset\frak H$ точки


$u_{0}$ уравнение $(2.2)$ может быть приведено  к виду


\begin{gather}


\Dot\xi{}_{q}=\xi_{q-1}+g_{q-1}(u),\quad q=1,2,\dotsc,p,\notag\\


0=\xi_{p}+g_{p}(u),\\


\Dot u{}^{1}=Su^{1}+F(u).\notag


\end{gather}


\end{thms}








\begin{pf}


В силу следствия 1.2 оператор


$M+N'_{u_{0}}$ \ $(L,\sigma)$-ограничен, причем $\infty$ --- полюс


порядка $p$ \ $L$-резольвенты оператора. Пусть


$\{\varphi_{1},\varphi_{2},\dotsc,\varphi_{p}\}$ --- цепочка


$(M+N'_{u_{0}})$-присоединенных векторов вектора


$\varphi_{0}\in\ker L\setminus\{0\}$. Нетрудно показать, что


векторы $\{\varphi_{0},\varphi_{1},\dotsc,\varphi_{p}\}$ линейно


независимы [6]. Поэтому любой вектор  $u\in\frak U$ может быть


представлен в виде


$$


u=u^{0}+u^{1}=\sum^{p}_{q=0}\xi_{q}\varphi_{q}+u^{1},


$$


а уравнение (2.2) приобретает вид (2.3), (2.4) с той лишь разницей,


что  $M_{k}$ есть сужение на $\frak U^{k}$ оператора


$M+N'_{u_{0}}$. Пусть


$\{\psi_{0},\psi_{1},\ldots,\psi_{p}\}\subset\frak U^{*}$ ---


биортогональный к $\{\varphi_{0},\varphi_{1},\dotsc,\varphi_{p}\}$


базис. Преобразовав (2.3) к виду


\begin{equation}


H\Dot u{}^{0}=u^{0}+M^{-1}_{0}(\Bbb I-Q)N(u),


\end{equation}


и подействовав на (2.8) последовательно функционалами


$\{\psi_{0},\psi_{1},\ldots,\psi_{p}\}$, с учетом (2.4) получим


(2.7).


\end{pf}





В ситуации теоремы 2.2 нельзя


говорить ни о существовании, ни о единственности решения задачи


(2.1), (2.2). Действительно, пусть (2.7) имеет вид


\begin{equation}


\Dot\xi{}_{2}=\xi_{1},\quad 0=\xi_{2}-\xi^{2}_{1}.


\end{equation}


Тогда задача Коши


\begin{equation}


(\xi_{1}(0),\xi_{2}(0))=(0,0)


\end{equation}


для системы (2.9), кроме стационарного


$(\xi_{1}(t),\xi_{2}(t))\equiv(0,0)$, имеет еще одно решение


$(\xi_{1}(t),\xi_{2}(t))\equiv(\frac{t}{2},\frac{t^{2}}{4})$. Если


же (2.7) имеет вид


\begin{equation}


\Dot\xi{}_{2}=\xi_{1}+1,\quad 0=\xi_{2}-\xi^{2}_{1},


\end{equation}


то задача (2.10), (2.11) вообще не имеет решения.








\section{Особенности фазового пространства}





Редуцируем задачу (0.1), (0.2) к задаче (2.1), (2.2). Для


этого положим $\frak U=L_{4}$ (все функциональные пространства


определены в области $\Omega$). Операторы $L$, $M$, $N$ определим


формулами


\begin{gather*}


\la Lu,v\ra=\int_{\Omega}(\lambda uv-u_{x_{k}}v_{x_{k}})dx \quad\forall


u,v\in \overset{\circ}{W}{}^{1}_{2}, \\


%


\la Mu,v\ra=\alpha\int_{\Omega}uv\,dx,\quad


\la N(u),v\ra=\beta\int_{\Omega}u^{3}v\,dx \quad\forall u,v\in  L^{4},


\end{gather*}


где $\la\cdot,\cdot\ra$ --- скалярное произведение в смысле $L_{2}$.


При $1\leq n\leq4$ вложение


$\overset{\circ}{W}{}^{1}_{2}\hookrightarrow L_{4}$


плотно и непрерывно, поэтому оператор $L\in \Cl(\frak U,\frak F)$,


где $\frak F=W^{-1}_{2}$. В [7] показано, что в силу вложения


$(L_{4})^{*2}\simeq L_{\frac{4}{3}}\hookrightarrow W^{-1}_{2}$


оператор $N\in C^{\infty}(\frak U;\frak F)$.








\begin{lems}


Пусть $1\leq n\leq4$, тогда при любых


$\lambda\in\Bbb R$, $\alpha\in\Bbb R\setminus\{0\}$ 


оператор $M$ \ $(L,\sigma)$-огра\-ничен, 


причем $\infty$ --- устранимая особая точка


$L$-резольвенты оператора $M$.


\end{lems}








{\bf Доказательство.} Пусть $\lambda\notin\sigma(\Delta)$, где


$\sigma(\Delta)$ --- спектр однородной задачи Дирихле в области


$\Omega$ для оператора Лапласа $\Delta$. Тогда существует оператор


$L^{-1}\in L(\frak F;\frak U)$. Очевидно, оператор


$L^{-1}M\in L(\frak U)$, отсюда оператор $(\lambda


L-M)^{-1}=L(\lambda\Bbb I-L^{-1}M)^{-1}\in L(\frak F;\frak U)$.





Пусть $\lambda\in\sigma(\Delta)$, тогда $\ker L=


\mathrm{span}\,\{\varphi_{k}: \lambda_{k}=\lambda\}$, $\im L=(\ker


L)^{\perp}$, где $\{\varphi_{k}\}$ --- ортонормированное (в смысле


$L_{2}$) семейство собственных функций однородной задачи Дирихле в


области $\Omega$ для оператора Лапласа $\Delta$, соответствующих


собственным значениям $\{\lambda_{k}\}=\sigma(L)$,


которые занумерованы по невозрастанию с учетом кратности.


(Ортогональность тоже в смысле $L_{2}$.) Поскольку


оператор $L$ фредгольмов, то для доказательства воспользуемся


следствием 1.1. Пусть $\varphi\in\ker L\setminus\{0 \}$, т.\,е.


$$


\varphi=\sum_{\lambda_{k}=\lambda}a_{k}\varphi_{k},\quad


\sum_{\lambda_{k}=\lambda}|a_{k}|>0.


$$


Поэтому


$$


M\varphi=\alpha\sum_{\lambda_{k}=\lambda}a_{k}\varphi_{k}\notin


\im L.\quad\square


$$





Ограничимся случаем $\dim\ker L=1$ и


построим проектор $\Bbb I -Q=\la\cdot,\psi\ra$, где $\psi\in\ker L$,


$\|\psi\|_{L_{2}}=1$, и множество


$$


\frak H=\{u\in\frak U:\la Mu+N(u),\psi\ra=0\}.


$$


Здесь $\la\cdot,\cdot\ra$ --- скалярное произведение в $L_{2}$.





Представив вектор $u$ в виде $u=s\psi+v$, где $v\in\frak


U^{1}=\{u\in\frak U:\la u,\psi\ra=0\}$, заметим, что 


множество $\frak H$ \ $C^{\infty}$-диффеоморфно множеству


$$


\frak H_{s}=\bigg\{(s,v)\in\Bbb R\times\frak U:


s^{3}\|\psi\|^{4}_{\frak U}+3s^{2}\int_{\Omega}\psi^{3}vdx+


s\bigg(3\int_{\Omega}\psi^{2}v^{2}dx+\alpha\beta^{-1}\bigg)+


\int_{\Omega}\psi v^{3}dx=0\bigg\}.


$$


В [5] множество $\frak H_{s}$ названо 2-сборкой Уитни, 


в [4] показано, что в случае $\alpha\beta>0$


при любом векторе $v\in\frak U^{1}$ существует точно одно число


$s\in\Bbb R$ такое, что $s\psi+v\in\frak H$. Доказательство этого


основано на том, что в любой точке $u\in\frak U$ вектор


$\psi\in\ker L\setminus\{0\}$ не имеет $(M+N'_u)$-присоединенных


векторов.





Считая, что $\alpha\beta<0$, выделим в множестве $\frak H$


точки складки, которые определяют множество


$$


\frak H'=\{u\in\frak


U: \la M\psi+N'_{u}\psi,\psi\ra =0\}\cap\frak H.


$$


Здесь $N'_{u}$ --- первая производная Фреше оператора $N$ в точке $u$,


$$


\la N'_{u}v,w\ra =3\beta\int_{\Omega}u^{2}vw\,dx \quad\forall


u,v\in L_{4}.


$$


Это множество $\frak H'$ \ $C^{\infty}$-диффеоморфно множеству


$$


\frak H'_{s}=\bigg\{(s,v)\in\Bbb R\times\frak U^{1}:


3s^{2}\|\psi\|^{4}_{\frak U}+6s\int_{\Omega}\psi^{3}v\,dx+


3\int_{\Omega}\psi^{2}v^{2}dx+\alpha\beta^{-1}=0\bigg\}\cap\frak H_{s}.


$$





В дальнейшем ограничимся случаем $\lambda=-\lambda_{1}$, где


$\lambda_{1}$ --- наибольшее собственное значение однородной задачи


Дирихле для оператора Лапласа $\Delta$ в области $\Omega$. Это


ограничение является частным случаем условия $\dim\ker L=1$.








\begin{thms}


Пусть $1\leq n\leq4$, $\alpha\beta<0$,


$\lambda=-\lambda_{1}$. Тогда


\begin{itemize}


\item[(i)] любой вектор


$\xi\in\ker L\setminus\{0\}$ не имеет


$(M+N'_{u})$-присоединенных векторов, если $u\in\frak H\setminus\frak H';$


\item[(ii)] любой вектор


$\xi\in\ker L\setminus\{0\}$ имеет точно один


$(M+N'_{u})$-присоединенный вектор, если $u\in\frak H'\setminus\frak U'$.


\end{itemize}


\end{thms}








\begin{pf}


(i) Пусть 


$$


u\in\frak H\setminus\frak H'=\{u\in\frak U:


\la Mu+N(u),\psi\ra =0,\ \la M\psi+N'_{u}\psi,\psi\ra\ne0\}.


$$


Любой вектор


$\xi\in\ker L\setminus\{0\}$ имеет вид $\xi=a\psi$, $a\in\Bbb


R\setminus\{0\}$. Поскольку $\la M\xi+N'_{u}\xi,\psi\ra\ne0$, то


$M\xi+N'_{u}\xi\notin\im L$.





(ii) Пусть $u\in\frak H'$. Тогда $M\psi+N'_{u}\psi=f\in


\im L$. Построим множество $\coim L=\{u\in\dom


L:\la u,\psi\ra =0\}$. Линеал $\coim L$, снабженный ``нормой


графика'' $\tri\cdot\tri=\|\cdot\|_{\frak U}+\|L\cdot\|_{\frak F}$,


является банаховым пространством. Обозначим через $L_{1}$ сужение


оператора $L$ на $\coim L$. В силу отрицательной


определенности оператора $L_{1}$ (т.\,е. $-\la L_{1}u,u\ra\geq


\const\|u\|^{2}_{L_{2}}$) и теоремы Банаха существует оператор


$L^{-1}_{1}\in L(\im L,\coim L)$. Отсюда


вытекает существование вектора $\varphi=L^{-1}_{1}f\in


\coim L$ такого, что $\la \varphi,\psi\ra =0$. Далее,


$$


\la M\varphi+N'_{u}\varphi,\psi\ra=\int_{\Omega}(\alpha+3\beta


u^{2})\varphi\psi \,dx=\la \varphi,M\psi+N'_{u}\psi\ra=


\la\varphi,f\ra=\la L^{-1}_{1}f,f\ra\leq0,


$$


причем равенство $\la L^{-1}_{1}f,f\ra=0$ имеет место тогда и только


тогда, когда $f=0$. Если же $\la L^{-1}_{1}f,f\ra<0$, то это значит, что


$M\varphi+N'_{u}\varphi\notin\im L$ и, следовательно, второго


$(M+N'_{u})$-при\-со\-е\-ди\-нен\-ного вектора вектор $\psi$ не имеет. Поэтому


и любой вектор $\xi=a\psi\in\ker L\setminus\{0\}$ обладает тем же


свойством.





Предположим, что $f=0$. Тогда $\la M\psi+N'_u\psi,w\ra=0$ при любом


$w\in\frak U$. Положим $w=u$, получим 


\begin{equation} 


3\int_\Omega u^3\psi \,dx+s\alpha\beta^{-1}=0,


\end{equation}


где $s$ из представления


$u=s\psi+v$, $v\in\frak U^{1}$. Поскольку $u\in\frak H$, то


\begin{equation}


0=\la Mu+N(u),\psi\ra=\int_{\Omega}u^{3}\psi\,


dx+s\alpha\beta^{-1}.


\end{equation}


Из (3.1), (3.2) вытекает, что


$s=0$ и $u=v\in\frak U^{1}$, что противоречит условию.


\end{pf}





Условие $M\psi+N'_{u}\psi=0$ означает,


что $\ker L\cap\ker(M+N'_{u})=\ker L\ne\{0\}$. Отсюда $\Bbb C$ ---


\mbox{$L$-спектр} оператора $M+N'_{u}$.


Поэтому оператор $M+N'_{u}$ не является $(L,\sigma)$-ограниченным.








\begin{thebibliography}{9}





\bibitem{1}


Hoff N.J. {\em Creep buckling} // Aeron. -- Quarterly 7. -- 1956.


-- \No\,1. -- P.\,1--20.


%%%%%%% ^^^^^^^^^^^^^^^^^^^^^^^^^^^^^^^^^^ verify via inet.


\bibitem{2}


Сидоров Н.А., Романова О.А. {\em О применении некоторых результатов


теории ветвления при решении дифференциальных уравнений} //


Дифференц. уравнения. -- 1983. -- Т.\,19. -- \No\,9. -- С.\,1516--1526.


\bibitem{3}


Свиридюк Г.А. {\em Фазовые пространства полулинейных уравнений типа


Соболева с относительно сильно секториальным оператором} // Алгебра


и анализ. -- 1994. -- Т.\,6. -- \No\,5. -- С.\,252--272.


\bibitem{4}


Свиридюк Г.А., Казак В.О. {\em Фазовое пространство 


начально-краевой задачи для уравнения Хоффа} // Матем. заметки. --


2002. -- Т.\,71. -- \No\,2. -- С.\,292--297.


\bibitem{5}


Бокарева Т.А., Свиридюк Г.А. {\em Сборки Уитни фазовых пространств


некоторых полулинейных уравнений типа Соболева}


// Матем. заметки. -- 1994. -- Т.\,55. -- \No\,3. -- С.\,3--10.


\bibitem{6}


Sviridyuk G.A., Fedorov V.E. {\em Sobolev type equations and


degenerate semigroups of operators}. -- Utrecht: VSP, 2003. -- 216 p.


\bibitem{7}


Свиридюк Г.А. {\em Квазистационарные траектории полулинейных


динамических уравнений типа Соболева} // Изв. РАН. Cер. матем.


-- 1993. -- Т.\,57. -- \No\,3. -- С.\,192--207.


\bibitem{8}


Ленг С. {\em Введение в теорию дифференцируемых многообразий}. --


Волгоград: Платон, 1996. -- 203~c.





\end{thebibliography}





\vskip 2cm





\postup{\it Челябинский государственный}{\qquad \it Поступили}


\postup{\it университет}{\it первый вариант $06.05.2003$}


\postup{\it Магнитогорский государственный}{\it окончательный вариант $26.03.2004$}


\postup{\it университет}{}





\label{end}


\end{document}





